\SourceLeaf{008}
\begin{center}
{\large THEORIA\par}
\medskip
{\Large\bfseries COMBINATIONIS OBSERVATIONUM\par}
\medskip
{\large ERRORIBUS MINIMIS OBNOXIAE.\par}
\medskip
{\large\textsc{Pars Prior.}\par}
\medskip
\rule{.22\linewidth}{.3pt}
\end{center}
\Article{1}
Quantacunque cura instituantur observationes, rerum naturalium magnitudinem spectantes, semper tamen erroribus maioribus minoribusve obnoxiae manent. Errores observationum plerumque non sunt simplices, sed e pluribus fontibus simul originem trahunt: horum fontium duas species probe distinguere oportet. Quaedam errorum caussae ita sunt comparatae, ut ipsarum effectus in qualibet observatione a circumstantiis variabilibus pendeat, inter quas et ipsam observationem nullus nexus essentialis concipitur: errores hinc oriundi irregulares seu fortuiti vocantur, quatenusque illae circumstantiae calculo subiici nequeunt, idem etiam de erroribus ipsis valet. Tales sunt errores ab imperfectione sensuum provenientes, nec non a caussis extraneis irregularibus, e. g. a motu tremulo aeris visum tantillum turbante: plura quoque vitia instrumentorum vel optimorum huc trahenda sunt, e. g. asperitas partis interioris libellularum, defectus firmitatis absolutae etc. Contra aliae errorum caussae in omnibus observationibus ad idem genus relatis natura sua effectum vel absolute constantem exserunt, vel saltem talem, cuius magnitudo secundum legem determinatam unice a circumstantiis, quae tamquam essentialiter cum observatione nexae spectantur, pendet. Huiusmodi errores constantes seu regulares appellantur.

Ceterum perspicuum est, hanc distinctionem quodammodo relativam esse, et a sensu latiore vel arctiore, quo notio observationum ad idem genus pertinentium accipitur, pendere. E. g. vitia irregularia in divisione instrumentorum ad \SourceLeaf{009}angulos mensurandos errorem constantem producunt, quoties tantummodo de observatione anguli determinati indefinite repetenda sermo est, siquidem hic semper eaedem divisiones vitiosae adhibentur: contra error ex illo fonte oriundus tamquam fortuitus spectari potest, quoties indefinite de angulis cuiusvis magnitudinis mensurandis agitur, siquidem tabula quantitatem erroris in singulis divisionibus exhibens non adest.

\Article{2}
Errorum regularium consideratio proprie ab instituto nostro excluditur. Scilicet observatoris est, omnes caussas, quae errores constantes producere valent, sedulo investigare, et vel amovere, vel saltem earum rationem et magnitudinem summo studio perscrutari, ut effectus in quavis observatione determinata assignari. adeoque haec ab illo liberari possit, quo pacto res eodem redit, ac si error omnino non affuisset. Longe vero diversa est ratio errorum irregularium, qui natura sua calculo subiici nequeunt. Hos itaque in observationibus quidem tolerare, sed eorum effectum in quantitates ex observationibus derivandas per scitam harum combinationem quantum fieri potest extenuare oportet. Cui argumento gravissimo sequentes disquisitiones dicatae sunt.

\Article{3}
Errores observationum ad idem genus pertinentium, qui a caussa simplici determinata oriuntur, per rei naturam certis \textit{limitibus} sunt circumscripti, quos sine dubio exacte assignare liceret, si indoles ipsius caussae \textit{penitus} esset perspecta. Pleraeque errorum fortuitorum caussae ita sunt comparatae, ut secundum legem continuitatis omnes errores intra istos limites comprehensi pro possibilibus haberi debeant, perfectaque caussae cognitio etiam doceret, utrum omnes hi errores aequali facilitate gaudeant an inaequali. et quanta probabilitas relativa, in casu posteriore, cuivis errori tribuenda sit. Eadem etiam respectu erroris totalis, e pluribus erroribus simplicibus conflati, valebunt, puta inclusus erit certis limitibus, (quorum alter aequalis erit aggregato omnium limitum superiorum partialium, alter aggregato omnium limitum inferiorum); omnes errores intra hos limites possibiles quidem erunt, sed prout quisque infinitis modis diversis ex erroribus partialibus componi potest, qui ipsi magis minusve probabiles sunt, alii maiorem, alii minorem facilitatem tribuere debebimus, eruique poterit lex probabilitatis re\SourceLeaf{010}%
lativae, si leges errorum simplicium cognitae supponuntur, salvis difficultatibus analyticis in colligendis omnibus combinationibus.

Exstant utique quaedam errorum caussae, quae errores non secundum legem continuitatis progredientes, sed discretos tantum, producere possunt, quales sunt errores divisionis instrumentorum, (siquidem illos erroribus fortuitis annumerare placet): divisionum enim multitudo in quovis instrumento determinato est finita. Manifesto autem, hoc non obstante, si modo non omnes errorum caussae errores discretos producant, complexus omnium errorum totalium possibilium constituet seriem secundum legem continuitatis progredientem, sive plures eiusmodi series interruptas, si forte, omnibus erroribus discretis possibilibus secundum magnitudinem ordinatis, una alterave differentia inter binos terminos proximos maior evadat, quam differentia inter limites errorum totalium, quatenus e solis erroribus continuis demanant. Sed in praxi casus posterior vix umquam locum habebit, nisi divisio vitiis crassioribus laboret.

\Article{4}
Designando facilitatem relativam erroris totalis \(x\), in determinato observationum genere, per characteristicam \(\varphi x\), hoc, propter errorum continuitatem, ita intelligendum erit, probabilitatem erroris inter limites infinite proximos \(x\) et \(x+\mathrm{d}x\) esse \(=\varphi x.\mathrm{d}x\). Vix, ac ne vix quidem, umquam in praxi possibile erit, hanc functionem a priori assignare: nihilominus plurá generalia eam spectantia stabilire possunt, quae deinceps proferemus. Obvium est. functionem \(\varphi x\) eatenus ad functiones discontinuas referendam esse, quod pro omnibus valoribus ipsius \(x\) extra limites errorum possibilium iacentibus esse debet \(=0\); intra hos limites vero ubique valorem positivum nanciscetur (omittendo casum, de quo in fine art. praec. locuti sumus). In plerisque casibus errores positivos et negativos eiusdem magnitudinis aeque faciles supponere licebit, quo pacto erit \(\varphi(-x)=\varphi x\). Porro quum errores leviores facilius committantur quam graviores, plerumque valor ipsius \(\varphi x\) erit maximus pro \(x=0\), continuoque decrescet, dum \(x\) augetur.

Generaliter autem valor integralis \(\int\varphi x.\mathrm{d}x\), ab \(x=a\) usque ad \(x=b\) extensi exprimet probabilitatem, quod error aliquis nondum cognitus iaceat inter limites \(a\) et \(b\). Valor itaque istius integralis a limite inferiore omnium errorum possibilium usque ad limitem superiorem semper erit \(=1\). Et quum \(\varphi x\)
